Scaling has driven most recent progress in language modelling, but it obscures a basic question: under a fixed data budget, which design choices actually buy generalization, and which only seem to until they are tested with matched controls? The BabyLM challenge~\cite{Warstadt2023, Hu2024BabyLM} makes this question tractable by holding the data fixed, so that what survives is informative about data-efficient pretraining. This paper, our BabyLM~2026 submission, asks whether linguistic structure from classical Pāṇinian grammar, refined over more than a millennium~\cite{Matilal1990WordWorld} and operationalized through morpheme-boundary embeddings and role-aware masking, improves small-model pretraining, and tests that claim rather than assuming it.

Pāṇinian grammar encodes compositional structure morphologically, through segmentation and inflection, and at the argument level, through kāraka categories that assign each participant a semantic role defined by its relation to the action rather than by surface order~\cite{KiparskyStaal1969, Cardona1974Karaka}. Because these roles are semantic, they are a plausible language-independent target for a pretraining signal even in English. Prabhāsa investigates this through a pre-registered hypothesis hierarchy with two mechanistically distinct levers. RQ-A (masking causality) asks whether, at matched mask-rate budget, stratifying token masking by kāraka role changes downstream syntax, realized via Vidyut morpheme-boundary N-hot embeddings~\cite{Prasad2024Vidyut}. RQ-B (auxiliary objective) asks whether a supervised token-level kāraka-role prediction head (a computational analogue of śābdabodha, verbal cognition) supplies a signal masking alone does not. Primary evaluation targets the BLiMP agreement and argument-structure subsets, with pre-registered thresholds of at least $+1.0$ point on the targeted subset, assessed across at least three seeds with paired bootstrap resampling and Holm--Bonferroni correction.

Our findings reorder the usual expectation. The clearest, most reproducible effect comes not from the Pāṇinian mechanisms but from the pretraining objective, and it is scale-dependent: a pure MLM objective is neutral against a hybrid MLM\,+\,CLM objective at the 10M-token budget but outperforms it by 5.49 BLiMP points at the 100M-token budget, placing our result in direct dialogue with the BabyLM~2024 winner GPT-BERT~\cite{CharpentierSamuel2024}, whose hybrid recipe we find is not uniformly preferable as scale grows. Against this, both Pāṇinian levers return pre-registered null or weak results once their confounds are removed (Figure~\ref{fig:forest}). These outcomes are a feature of the methodology, not a disappointment of it: an earlier single-seed run had suggested a $+1.23$-point gain from kāraka masking that the matched-budget comparison shows was a mask-rate artifact. The contribution is therefore a rigorous separation of the levers that matter (objective and architecture) from a linguistically motivated prior that does not at this scale, with all code, data, and per-seed results released openly.
